\documentclass[10pt]{article}

\usepackage[utf8]{inputenc}

\usepackage[T1]{fontenc}

\usepackage{amsmath}

\usepackage{amsfonts}

\usepackage{amssymb}

\usepackage[version=4]{mhchem}

\usepackage{stmaryrd}

\usepackage{bbold}

\usepackage{mathrsfs}

\usepackage{graphicx}

\usepackage[export]{adjustbox}

\graphicspath{ {./images/} }



\DeclareUnicodeCharacter{0131}{$\imath$}



\begin{document}

ТРЕУГОЛЬНОЕ ЧИСЛО - см. А рифметический ряд. ТРЕУГОЛЬНЫИ МЕТОД СУММИРОВАНИЯ матричный метод суммирования, определенный треугольной матрицей



\[

A=\left\|a_{n k}\right\|, n, k=1,2, \ldots,

\]



т. с. матрицей, у к-рой $a_{n k}=0$ при $k>n$. Т. м. с. является частным случаем конечиострочных методов сумжирования. Треугольная матрица $A$ наз. п о р м а л ьн о й, если $a_{n n} \neq 0$ для всех $n$. Преобразованше



\[

\sigma_{n}=\sum_{k=1}^{n} a_{n k} s_{k}

\]



осуществляөмое посрепстьом нормальной треугольноӥ матрицы $A$, допускает обращение



\[

s_{n}=\sum_{k=1}^{n} a_{n k}^{-1} \sigma_{k}

\]



где $A^{-1}=\left\|a_{n k}^{-1}\right\|$ - матрица, обратная для $A$. Этот факт упрощает доказательство ряда теорем для матричных методов сумми рования, определенных нормальными треугольными матрицами. 1 Т Т. М. с. относятся, напр., Чезаро метод суммирования, Воропого метод суммирования.



JIит. [1] Х а р д и Г., Расходящиеся ряды, пер. с англ., M., 1951, [2] К у к P., Бесконечные матрицы и пространства последоватсльностей, пер. с англ., М., 1960; [3] Б а p о н С., Введение в теорию суммируемости рядов, Таллин, 1977.



ТРЕУГОЛЬНЫИ ЭЛЕМЕН эндоморфизмов конечномерного векторного пространства $V$ над полем $k$ - элемент $X \in$ End $V$, все собствепные значепия к-рого принадлежат $k$. Если $k$ алгебраически замкнуто, то всякий элемент из End $V$ треуголен. Для Т. э. $X$ (и только для такого элемента) существует базис в $V$, относительно к-рого матрица эндоморфизма $X$ треугольна (или, что то же, существует инвариантный относительно $X$ полный флаг в $V$ ). Для Т. ә. имеется Жо ордана разложение над $k$. Существует ряд обобщений понятия Т. э. в End $V$ на случай бесконсчномерного $V$ (cM. [2]).



\begin{enumerate}

  \setcounter{enumi}{1}

  \item Т. э. конечномерной алгебры $A$ над полем $k-$ такой элемент $a \in A$, что оператор правого (или левого, в зависимости от рассматриваемого случая) умножения на $a$ является T. э. в алгебре $\operatorname{End}_{k} A$. Если $A$ изоморфна алгебре End $V$ для нек-рого конечномерного векторного пространства $V$ над $k$, то эти два (формально различные) ошределения приводят к одному и тому же понятию.

\end{enumerate}



В алгебрах Ли треугольность элемента $x \in A$ означает треугольность эндоморфизма $\operatorname{ad}_{x}$ (где $\left.\operatorname{ad}_{x}(y)=[x, y]\right)$. Множество всех Т. э. в алгебре Ли, вообще говоря, не замкнуто относительно онераций сложения и коммутирования (напр., для $\& l(2, \mathbb{R})$ - простой алгебры Ли вещественных матриц порядка 2 со следом 0). Однако в случае разрешимой алгебры $A$ әто множество является даже характеристич. идеалом в $A$.



\begin{enumerate}

  \setcounter{enumi}{2}

  \item Т. э. в связной конечномерной группе Ли $G-$ элемент $g \in G$, для к-рого $\mathrm{Ad}_{g}$ является Т. э. в End $g$ (здесь Ad : $G \rightarrow$ Aut $g$ - присоеддиенное представление группы Ли $G$ в групге Aut $\mathfrak{g} \subset$ End $q$ автоморфизмов ее алгебры Ји (̧). Если $\exp : \mathfrak{g} \rightarrow \ddot{G}-$ экспоненциальное отображение, а $X \in \mathfrak{T}$ - Т. ә. (в смысле пункта 2), то $\exp (X)-$ Т. ә. в $\mathscr{G}$. Обратное утверждение в общем слутае певершо.

\end{enumerate}



Алгебры Ли и группы Ли, все элементы к-рых треугольны, наз. треугольными алгебрами илш группами соответственно, а также Ли вполие разрешилыми алгебрами и Ли вполне разрешимыли группами.



Juт. [1] Б о р е ль А., Линейные алгебраические групшы, пер. с англ., М., 1972, [2] П л о т к и н Б. И., Группы автоморфизмов алгебрапческих систем, М., 1966; [3] II ос тн ик о в М. М., Јинейная алгебра и дифференциальная геометрия, М., 1979.

B. В. Горбачевич. ТРЕФФЦА МЕТОД - один из вариационных методов решения краевых задач. Іусть требуется решить краевую задачу



\[

\Delta u=0,\left.\quad u\right|_{s}=\varphi,

\]



где $S$ - граница области $\Omega \subset \mathbb{R}^{m}$. Репение задачи (*) минимизи рует функционал



\[

J(u)=\int_{\Omega}(\operatorname{grad} u(x))^{2} d x

\]



среди всех функциі̆, удовлетворяющих краевому условию $u / S=\varphi$. Т. м. заключается в следующем. Пусть дана последовательность гармонических в $\Omega$ функций $v_{1}, v_{2}, \ldots$, квадратично суммируемых в $\Omega$ вмес'те с первыми гроизводными. Приближенное решение разыскивается в виде



\[

u_{n}=\sum_{j=1}^{n} c_{j} v_{j}

\]



коэффициенты $c_{j}$ определяются из условия минимума $J\left(u_{n}-u\right)$, где $u$ - точное решение задачи (*). Это приводит к следующей системе уравнений относительно $c_{1}, \ldots, c_{n}$



\[

\begin{gathered}

\sum_{j=1}^{n} c_{j} \int_{S} v_{j} \frac{\partial v_{i}}{\partial v} d S=\int_{S} \varphi \frac{\partial v_{i}}{\partial v} d S, \\

i=1, \ldots, n,

\end{gathered}

\]



где $v$ - внешняя нормаль к $S$.



Т. м. допускает обобщение и обоснование для различных краевых задач (см. [2] - [4])



Метод предложен Э. Трефффцем (см. [1]).



JIum.. [1] T r e f f t z E., в cб.: Verhandl, des 2. Internat. Kongresses für technische Mechanik. Zürich, 1926, 12-17 Sept., Zürich - Lpz., 1927, S. 131-37; [2] M и х лі и н С. Г., Вариапионные методы в математической физике, 2 изд., іу., 1970;



\includegraphics[max width=\textwidth, center]{2023_07_09_c0584486ec4a24b637acg-1}

ны й П. И., Вычислительные методы высшей математики, т. 2, механ. и астр.», 1956, Nํ1 13, в. 3, с. 69-89. Г. М. Вайıикюо.



«TPEХ СИГМ ПРАВИЛО»- мнемоническое правило, согласно к-раму в нек-рых задачах теории вероятностей и математич. статистики считают практически невозможным событие, заключающееся в отклонении значения нормально распределенной случайной величины от ее математич. ожидания больпе, чем на три стандартных отклонения.



Пусть $X$ - нормально $N\left(a, \sigma^{2}\right)$ распределенная слуqайная величина. Для лгобого $k>0$



\[

\mathrm{P}\{|X-a|<k \sigma\}=2 \Phi(k)-1,

\]



гдо Ф (•) - функция распределения стандартного нормального закона, откуда, в частности, при $k=3$ следует, чтO



\[

\mathrm{P}\{a-3 \sigma<X<a+3 \sigma\}=0,99730 .

\]



Последнее равенство означает, что значения случайной величины $X$ будут отклоняться от ее математич. ожиданил а на расстояние, превосходящее $3 \sigma$ в среднем не более чем 3 раза на одну тысячу испытаний. Именно это обстоятельство используют иногда экспериментаторы в нек-рых задачах теории вероятностей и математич. статистики, считая, что событие $\{|X-a|>$ $>3 \sigma\}$ является практически невозможным и, следовательно, событие $\{|X-a|<3 \sigma\}-$ практически достоверным. В этом случае говорят, что экспериментатор придеряиивается «правила трех сигм».



Jum.: [1] С м и р н о в Н. В., Д у н и н-Б а с к и й И. В., Курс теории вероятностей и математической статистики для технических приложений, 3 изд., М., 1969.



ТРЕХ ТЕЛ ЗАДАЧА - задача 0 двияении Т тел, рассматриваемых как материальные точки, взаимно притягиваюпцхся поо закону тяготенія Ньютона. Классич. пример Т. т. з.- задача о движении системы Солнце - Земля - Јуна. Т. т. з. состоит в нахожде-





\end{document}